\section{Bubble Sort: an Invariant Indexed by the Iteration}\label{sec:bs}

\subsection{The program}
\begin{lstlisting}
method main(a: array<int, int>, n: int) returns (b: array<int, int>)
  requires 0 <= n
  ensures sorted(b, 0, n)
  ensures perm(b, a)
{
  var i: int; var j: int; var t: int;
  b := a; i := n;
  while 1 < i {
    j := 0;
    while j + 1 < i {
      if b[j + 1] < b[j] { t := b[j]; b[j] := b[j + 1]; b[j + 1] := t; }
      j := j + 1;
    }
    i := i - 1;
  }
}
\end{lstlisting}

\subsection{The trace equation system}
The outer loop modifies only $b$, and only through the inner loop, so its trace is the chain of the inner
results, which the fold \lemma{results} names:
\begin{align}
  K &= \cnt(0, 1, n') \qquad I = \cnt(n, -1, n') \qquad b = \cur(B) \label{eq:bs0}\\
  B &= \mathsf{results}(a, \bl{K}, \seq{B}) \quad\text{where}\quad \mathsf{results}(x_0, ks, bs) = \tr{x_0}{ks, bs}{x, k, b \mapsto b} \nonumber\\
  &\al{\bl{K}, \bl{I}, \seq{B}, \seq{B}'}{k, i, b, B' \mapsto b = \cur(B')} \label{eq:bs1}\\
  &\al{\bl{K}, \bl{B}, \bl{I}, \seq{B}', \seq{J}'}{k, x, i, B', J \mapsto B' = \mathsf{bubble\_pass}(x, \bl{J})} \label{eq:bs2}\\
  &\qquad\text{where}\quad \mathsf{bubble\_pass}(b, js) = \tr{b}{js}{x, j \mapsto {}\nonumber\\
  &\qquad\qquad\qquad\qquad \upd{x[j+1] < x[j]}{\swap{x}{j}{j + 1}}\ \ \upd{x[j] \le x[j+1]}{x}} \nonumber\\
  &\al{\bl{K}, \bl{I}, \seq{J}}{k, i, j \mapsto n \le j + k + 1} \label{eq:bs3}\\
  &\al{\bl{K}, \bl{I}, \seq{J}'}{k, i, J \mapsto \al{\bl{J}}{j \mapsto j + k + 1 < n}} \label{eq:bs4}\\
  &\al{\bl{K}, \bl{I}, \seq{J}, \seq{J}'}{k, i, j, J \mapsto j = \cur(J)} \label{eq:bs5}\\
  &\al{\bl{K}, \bl{I}, \seq{K}', \seq{J}'}{k, i, K', J \mapsto J = \sums(0, \mapc{\bl{K'}}{1})} \label{eq:bs6}\\
  &\ldots \wedge\ n \le \intf(n') + 1\ \Longrightarrow\ \sorted(b, 0, n) \wedge \perm(b, a) \nonumber
\end{align}
The inner trace folds into $\mathsf{bubble\_pass}(x, \bl{J})$, a pass over the positions of the inner
counter $J$, by a \texttt{\{:fold\}} function of the file, like $\sums$ in \texttt{teq.dfy}. The test and
the exit are stated with $i = n - k$ already substituted: the pass of iteration $k$ runs over the
positions $0, \ldots, n - k - 2$.

\subsection{The summary}
\begin{align*}
  &\lemma{current\_results\_bubble}:\\
  &\quad \al{ks, bs, Bs}{k, b, B' \mapsto b = \cur(B')}\\
  &\quad\wedge\ \al{ks, \bl{\mathsf{results}(a, ks, bs)}, Bs, Js}{k, x, B', J \mapsto B' = \mathsf{bubble\_pass}(x, \bl{J})}\\
  &\quad\wedge\ \al{ks, js}{k, j \mapsto len \le j + k + 1}\\
  &\quad\wedge\ \al{ks, Js_2}{k, J \mapsto \al{\bl{J}}{j \mapsto j + k + 1 < len}}\ \wedge\ Js_2 = Js\\
  &\quad\wedge\ \al{ks, js_3, Js_3}{k, j, J \mapsto j = \cur(J)}\ \wedge\ js_3 = js\ \wedge\ Js_3 = Js\\
  &\quad\wedge\ \al{ks, Ks, Js}{k, K', J \mapsto J = \sums(0, \mapc{\bl{K'}}{1})}\\
  &\quad\Longrightarrow\ \cur(\mathsf{results}(a, ks, bs))\\
  &\qquad\qquad = \cur(\tr{a}{ks}{b, k \mapsto \cur(\mathsf{bubble\_pass}(b, \bl{\cnt(0, 1, \natf(len - k - 1))}))})
\end{align*}
This is the summary in its original form (October 9), before the fragment lemmas: the inner counter is
matched as $\sums$, and the test and exit bind $len$. The left-hand side $\cur(\mathsf{results}(a, ks, bs))$
is matched against $b = \cur(B)$ through \eqref{eq:bs0}, binding $ks := \bl{K}$ and $bs := \seq{B}$;
\eqref{eq:bs1} binds $Bs := \seq{B}'$, \eqref{eq:bs2} binds $Js := \seq{J}'$, \eqref{eq:bs3} binds
$js := \seq{J}$ and $len := n$; \eqref{eq:bs4}, \eqref{eq:bs5}, \eqref{eq:bs6} are matched through
duplicates and the equations identify them, as in Section~\ref{sec:sel:sum}. The result is the trace of passes, the
$k$-th over $\bl{\cnt(0, 1, \natf(n - k - 1))}$. Note that the arithmetic $n - k - 1$ on the bound
variable $k$ is admissible on the right-hand side of a rule; on a left-hand side it would not be matched,
which is why the original version of this file had a function $\mathsf{passlen}(len, k)$ for it, which
the induction below makes unnecessary.

\subsection{The induction with a clamped invariant}
After $k$ passes the suffix $b[n - k..n)$ is sorted and all of $b[0..n - k)$ are at most all of it. Stated
with $\sorted$ and $\partitioned$ as in Section~\ref{sec:sel:ind}, the invariant makes sense only for $k \le n$, but
the number of outer iterations $n'$ is known only as $n \le n' + 1$ together with the iteration
constraint $1 < i$, a disjunction $n' = 0 \vee (2 < n \wedge n' < n)$ that the simplifier does not split,
so the induction cannot fix the number of steps to $n - 1$ (as Section~\ref{sec:sel:ind} fixes it to $n$). Instead the
invariant is a predicate indexed by the iteration, $\mathsf{bubbled}(b, k, n)$, which means the sorted
suffix of length $k$ once $k \le n$ and the whole range sorted once $k \ge n$; then every step keeps it,
also the empty passes beyond $n - 1$:
\begin{align*}
  &\lemma{bubbled\_pass}:\\
  &\quad p = len - k - 1 \wedge \mathsf{bubbled}(b, k, len)\\
  &\quad\Longrightarrow \mathsf{bubbled}(\cur(\mathsf{bubble\_pass}(b, \bl{\cnt(0, 1, \natf(p))})), k + 1, len)
    \Longleftrightarrow \mathsf{true}\\[4pt]
  &\lemma{sorted\_bubbled\_steps}:\\
  &\quad len \le \intf(m) + 1\\
  &\quad\wedge\ \big(\forall b, k.\ 0 \le k < \intf(m) \wedge \mathsf{bubbled}(b, k, len) \Longrightarrow \mathsf{bubbled}(F(b, k), k + 1, len)\big)\\
  &\quad\Longrightarrow\ \sorted(\cur(\tr{a}{\bl{\cnt(0, 1, m)}}{b, k \mapsto F(b, k)}), 0, len)\ \Longleftrightarrow\ \mathsf{true}
\end{align*}
The induction matches with $F := \lambda b\, k.\ \cur(\mathsf{bubble\_pass}(b, \bl{\cnt(0, 1, \natf(n - k - 1))}))$,
$m := n'$, $len := n$. Its quantified condition is decided by \lemma{bubbled\_pass}, whose length $p$ is a
variable with the equation $p = len - k - 1$ as a side condition, since a pattern $len - k - 1$ with two
variables is not matched; the equation is closed after $p$ is bound to $n - k - 1$ by the left-hand side
and decided by arithmetic. The lemma is sound for any number of steps $m$ with $len \le m + 1$: the
invariant at $k = 0$ holds trivially, and $\mathsf{bubbled}(b, m, len)$ with $len - m \le 1$ is
$\sorted(b, 0, len)$. For $\perm$, \lemma{perm\_steps} (the quantified condition $\forall b, k.\ \perm(F(b, k), b)$,
decided by \lemma{perm\_bubble\_pass}) and \lemma{perm\_refl}.

\subsection{The idea}
Two things differ from selection sort. The inner loop is summarized to a function of the previous
array that is itself a trace, $\cur(\mathsf{bubble\_pass}(\ldots))$, so the outer trace is a trace of
traces; the lemma database needs no knowledge of what a pass does beyond the two facts
\lemma{bubbled\_pass} and \lemma{perm\_bubble\_pass}. And the invariant is written as one predicate over
the iteration, which absorbs the case distinction on the number of passes that the goal does not provide.
The pattern "invariant as a predicate of $(b, k)$, step lemma, generic induction" is the same as
\lemma{current\_trace\_bounded} of \texttt{teq.dfy}, with the induction hypothesis as an abstract atom.
